\documentclass{cys}
\usepackage[utf8]{inputenc}
\usepackage[T1]{fontenc}
\usepackage{geometry}
\usepackage{graphicx}
\usepackage{float}
\usepackage{hyperref}
\usepackage[english,indonesian]{babel}

\title{Optimasi Pemrosesan Template LaTeX Menggunakan Pendekatan Berbasis AI}
\author{Academic Support Agent}
\date{\today}

\begin{document}
\maketitle

\begin{abstract}

\end{abstract}

\chapter{PENDAHULUAN}


\chapter{TINJAUAN PUSTAKA}


% Original content (commented out):
% \documentclass{cys}
\usepackage[ansinew]{inputenc}

\usepackage{graphicx}



\title{Paper Title Here \\ Several Lines, Please Use Capital Letters (for First Letters) in English}

\author{John Smith$^1$, Juan Perez$^2$, and Ivan Petrov$^2$}

\affil{ 
$^1$Institute, City, \authorcr
Country             
\authorcr \authorcr
$^2$Institute, City, \authorcr
Country             
\authorcr  \authorcr
author1@xxx.xx, author2@xxx.xx
\authorcr  \authorcr
}


\begin{document}

\maketitle

\begin{abstract}
This paper presents...
\end{abstract}

\begin{keywords} 
Word1, word2, word3.
\end{keywords} 

% If you do not speak Spanish, we will translate the Spanish parts for you.
\othertitle{T\'{\i}tulo en espa\~nol, solo primera letra es may\'uscula} 
 
\begin{resumen} 
Resumen en espa\~nol.
\end{resumen} 

\begin{palabrasclave} 
Palabra1, palabra2, palabra3.
\end{palabrasclave} 

\section{Introduction}
\label{sec:introduction}
Some text with reference \cite{hulth:2003}.

More text with other reference \cite{rose:2010}.


\section{Related Work}
\label{sec:relatedWork}

Other text.

And more.

\begin{figure}[h]
\centering
\includegraphics[width=0.1\textwidth]{foto.JPG}
\caption{Snapshot of something}
\label{jdm1}
\end{figure}

More text.

\begin{table}[h]
	\centering
	\caption{Tables, no capital letters, no point at the end}
	\begin{tabular}{lrrrr}
	\hline
		Classes & Articles \\
	\hline
		one \emph{and} two& 61 & 115,729 \\
		three & 169 & 869,607 \\
		four & 209 & 2,402,704 \\
		five & 206 & 2,561,115 \\		
	\hline
	\end{tabular}
	\label{table:Table1}
\end{table}

\subsection{Linguistic Features}
\label{subsection:linguistic}

Features. 

\subsubsection{Lexical Features}
\label{subsection:lexicalFeatures}

Other features.

\begin{equation}
H(X|Y) = - \sum_{y \in Y} P(y_i) \sum_{x \in X} p(x_i|y_i)\log p(x_i|y_i)
\label{equation:conditionalEntropy1}
\end{equation}


\section{Conclusion and Future Work}
\label{sec:conclusionAndFutureWork}

Conclusions here.

\section*{Acknowledgements} 
We would like to thank.. 

This work is funded by...


% OBLIGATORY: use BIBTEX formatting!
\small{
\bibliographystyle{cys}
\bibliography{biblio}
}
\normalsize


\begin{biography}[foto.JPG]{John Smith} 
is professor ...
\end{biography}

\begin{biography}[foto.JPG]{Juan Perez} 
is researcher...
\end{biography}

\begin{biography}[foto.JPG]{Ivan Petrov} 
is researcher...
\end{biography}


{\vskip 12pt}
\noindent
\footnotesize {Article received on 06/12/2012; accepted on 16/01/2013.}

\end{document}


\end{document}